\scanpage{30}
\begin{center}$-\ 31\ -$\end{center}
\hfill n$^\circ$ 371

PROPOSITION 1.--- \uline{Soient} $E,F$ \uline{deux} $A$-\uline{algèbres}, $a$ (resp. $b$) \uline{un idéal bilatère de} $E$ (resp. $F$). \uline{Alors l'isomorphisme canonique}
\[
(E/a)\otimes(F/b)\longrightarrow(E\otimes F)/(\operatorname{Im}(a\otimes F)+\operatorname{Im}(E\otimes b))
\]
(chap. II, $3^{\mathrm e}$ éd., \S\ 3, n$^\circ$ 6, cor. 1 de la prop. 6) \uline{est un isomorphisme d'algèbres}.

COROLLAIRE 1.--- \uline{Soient} $E$ \uline{une} $A$-\uline{algèbre}, $a$ \uline{un idéal de} $A$. \uline{Alors le} $A$-\uline{module} $aE$ \uline{est un idéal bilatère de} $E$, \uline{et l'isomorphisme canonique de} $A$-\uline{modules}
\[
(A/a)\otimes_A E\longrightarrow E/aE
\]
\uline{est un isomorphisme de} $A$-\uline{algèbres}.

COROLLAIRE 2.--- \uline{Si} $a,b$ \uline{sont deux idéaux de} $A$, \uline{la} $A$-\uline{algèbre} $(A/a)\otimes(A/b)$ \uline{est canoniquement isomorphe à} $A/(a+b)$.

COROLLAIRE 3.--- \uline{Soient} $E,F$ \uline{deux} $A$-\uline{algèbres}, $a$ \uline{un idéal de} $A$ \uline{contenu dans l'annulateur de} $F$. \uline{Alors la} $(A/a)$-\uline{algèbre} $E\otimes_A F$ \uline{est canoniquement isomorphe à} $(E/aE)\otimes_{A/a}F$.

PROPOSITION 2.--- \uline{Si l'algèbre} $E$ (resp. $F$) \uline{est composée directe d'une famille finie d'idéaux bilatères} $(a_i)_{1\le i\le m}$ (resp. $(b_j)_{1\le j\le n}$), \uline{alors les applications canoniques} $a_i\otimes b_j\to\operatorname{Im}(a_i\otimes b_j)$ \uline{sont bijectives, et} $E\otimes F$ \uline{est composée directe des idéaux bilatères} $\operatorname{Im}(a_i\otimes b_j)$.

Cela résulte aussitôt du chap. II, $3^{\mathrm e}$ éd., \S\ 3, n$^\circ$ 7, prop. 7, et de la définition de la multiplication dans $E\otimes F$.

PROPOSITION 3.--- \uline{Soient} $A,B$ \uline{deux anneaux commutatifs}, $\rho:A\to B$ \uline{un homomorphisme d'anneaux}, $E,F$ \uline{deux} $A$-\uline{algèbres}. \uline{Alors l'isomorphisme canonique}
\[
\rho^*(E)\otimes_B\rho^*(F)\longrightarrow\rho^*(E\otimes_A F)
\]
(chap. II, $3^{\mathrm e}$ éd., \S\ 4, n$^\circ$ 1, prop. 3) \uline{est un isomorphisme de} $B$-\uline{algèbres}.

PROPOSITION 4.--- \uline{Soient} $A,B$ \uline{deux anneaux commutatifs}, $\rho:A\to B$ \uline{un homomorphisme d'anneaux}, $E$ \uline{une} $A$-\uline{algèbre}, $F$ \uline{une} $B$-\uline{algèbre}. \uline{\unsure{Alors} l'isomorphisme canonique}
\mspage{29}
